\documentclass[12pt]{article}
\usepackage[utf8]{inputenc}
\usepackage[T1]{fontenc}
\usepackage{hyperref}
\title{Tor\'a Language Resources for K-12 Curriculum Development}
\author{Compiled by Claude Sonnet 5 (claude-sonnet-5) for Andr\'as Kornai}
\date{}
\begin{document}
\maketitle

\section{Introduction}
Tor\'a (ISO 639-3 trz; Glottocode tora1263) is a Chapacuran language of the lower Marmelos River, Amazonas State, Brazil. Search results describe it as dormant or extinct. Sources agree that at most a few semi-speakers remained by 2018. Some sources say the Tor\'a people are interested in reviving it. Very little material exists, so this report is short.

\section{Primary Educational Resources}
No primers, textbooks or school programmes for Tor\'a turned up in the searches.

\section{Linguistic and Reference Materials}
The language is known mainly from a wordlist of roughly 670 words, phrases and sentences, collected by Curt Nimuend\-aj\'u and Euclydes Henrique do Valle Bentes in the 1920s.
\begin{itemize}
\item Tor\'a language (Wikipedia): \url{https://en.wikipedia.org/wiki/Tor%C3%A1_language}
\item Tor\'a document at the DICE project, University of Missouri (a PDF, seen in search results only; its content was not checked): \url{https://dice.missouri.edu/assets/docs/south-america-other/Tora.pdf}
\item Tor\'a numerals page, Max Planck Institute for Evolutionary Anthropology: \url{https://lingweb.eva.mpg.de/channumerals/Tora.htm}
\item Chapacuran languages (Wikipedia): \url{https://wikipedia.com/wiki/Chapacuran}
\end{itemize}

\section{Religious and Community Resources}
No religious or community resources were found.

\section{Digital and Audio Resources}
\begin{itemize}
\item Endangered Languages Project entry: \url{https://endangeredlanguages.com/lang/2140}
\item Verbix endangered-languages entry: \url{https://docs.verbix.com/EndangeredLanguages/Tora}
\end{itemize}
No audio or video was found.

\section{Curriculum Development Context}
The documentation consists of a 1920s wordlist, so any revitalization curriculum would have to be built from that wordlist and from comparison with related Chapacuran languages. The ethnic population is roughly 300 people (Wikipedia gives about 330 in 2014; one mirror gives 310 in 2012). Contact with the community and with Brazilian linguists would be needed. No such contacts were identified in the searches.
\end{document}